\documentclass[11pt]{article}
% Shared preamble of the RMC kernel write-ups (% Shared preamble of the RMC kernel write-ups (% Shared preamble of the RMC kernel write-ups (\input{preamble} after \documentclass).
\usepackage[margin=1in]{geometry}
\usepackage{amsmath,amssymb,bm}
\usepackage{graphicx}
\usepackage{booktabs}
\usepackage{hyperref}

\newcommand{\dd}{\,\mathrm{d}}
\newcommand{\Ein}{E}
\newcommand{\Eout}{E'}
\newcommand{\Ecm}{E_{c}}
\newcommand{\mucm}{\mu_{c}}
\newcommand{\mulab}{\mu_0}

% Common closing section: how the verification figures are produced
\newcommand{\verificationmethod}{%
Each figure is produced by \texttt{verification/verify\_laws.py}. For an incident
energy $\Ein$ along $+z$, $10^6$ collisions are sampled with MC/DC's own reaction
sampler (\texttt{sample\_elastic\_scattering}, \texttt{sample\_inelastic\_scattering}
or \texttt{sample\_fission}), and every emitted neutron is histogrammed in
$48 \times 24$ bins of $(\Eout, \mulab)$, with $\mulab$ the lab cosine to $+z$. The
inverted kernel used by RMC is integrated over the same bins (Gauss--Legendre panels;
two-body lines are split where $\mulab^*(\Eout)$ crosses a cosine edge), multiplied by
the number of collisions, and compared with a $\chi^2$ test over bins with at least 5
expected counts (the rest pooled into one bin). The tables list $\chi^2/\mathrm{dof}$,
its $p$-value, the emitted neutrons per collision (sampled vs.\ the integral of the
inverted kernel), and the largest relative change of a bin integral when the
quadrature panels are doubled. Left panels: outgoing-energy marginal; middle:
lab-cosine marginal; right: per-bin $z = (\text{sampled} - \text{inverted}) /
\sqrt{\text{inverted}}$.}
 after \documentclass).
\usepackage[margin=1in]{geometry}
\usepackage{amsmath,amssymb,bm}
\usepackage{graphicx}
\usepackage{booktabs}
\usepackage{hyperref}

\newcommand{\dd}{\,\mathrm{d}}
\newcommand{\Ein}{E}
\newcommand{\Eout}{E'}
\newcommand{\Ecm}{E_{c}}
\newcommand{\mucm}{\mu_{c}}
\newcommand{\mulab}{\mu_0}

% Common closing section: how the verification figures are produced
\newcommand{\verificationmethod}{%
Each figure is produced by \texttt{verification/verify\_laws.py}. For an incident
energy $\Ein$ along $+z$, $10^6$ collisions are sampled with MC/DC's own reaction
sampler (\texttt{sample\_elastic\_scattering}, \texttt{sample\_inelastic\_scattering}
or \texttt{sample\_fission}), and every emitted neutron is histogrammed in
$48 \times 24$ bins of $(\Eout, \mulab)$, with $\mulab$ the lab cosine to $+z$. The
inverted kernel used by RMC is integrated over the same bins (Gauss--Legendre panels;
two-body lines are split where $\mulab^*(\Eout)$ crosses a cosine edge), multiplied by
the number of collisions, and compared with a $\chi^2$ test over bins with at least 5
expected counts (the rest pooled into one bin). The tables list $\chi^2/\mathrm{dof}$,
its $p$-value, the emitted neutrons per collision (sampled vs.\ the integral of the
inverted kernel), and the largest relative change of a bin integral when the
quadrature panels are doubled. Left panels: outgoing-energy marginal; middle:
lab-cosine marginal; right: per-bin $z = (\text{sampled} - \text{inverted}) /
\sqrt{\text{inverted}}$.}
 after \documentclass).
\usepackage[margin=1in]{geometry}
\usepackage{amsmath,amssymb,bm}
\usepackage{graphicx}
\usepackage{booktabs}
\usepackage{hyperref}

\newcommand{\dd}{\,\mathrm{d}}
\newcommand{\Ein}{E}
\newcommand{\Eout}{E'}
\newcommand{\Ecm}{E_{c}}
\newcommand{\mucm}{\mu_{c}}
\newcommand{\mulab}{\mu_0}

% Common closing section: how the verification figures are produced
\newcommand{\verificationmethod}{%
Each figure is produced by \texttt{verification/verify\_laws.py}. For an incident
energy $\Ein$ along $+z$, $10^6$ collisions are sampled with MC/DC's own reaction
sampler (\texttt{sample\_elastic\_scattering}, \texttt{sample\_inelastic\_scattering}
or \texttt{sample\_fission}), and every emitted neutron is histogrammed in
$48 \times 24$ bins of $(\Eout, \mulab)$, with $\mulab$ the lab cosine to $+z$. The
inverted kernel used by RMC is integrated over the same bins (Gauss--Legendre panels;
two-body lines are split where $\mulab^*(\Eout)$ crosses a cosine edge), multiplied by
the number of collisions, and compared with a $\chi^2$ test over bins with at least 5
expected counts (the rest pooled into one bin). The tables list $\chi^2/\mathrm{dof}$,
its $p$-value, the emitted neutrons per collision (sampled vs.\ the integral of the
inverted kernel), and the largest relative change of a bin integral when the
quadrature panels are doubled. Left panels: outgoing-energy marginal; middle:
lab-cosine marginal; right: per-bin $z = (\text{sampled} - \text{inverted}) /
\sqrt{\text{inverted}}$.}

\usepackage{amsthm}

\newtheorem{proposition}{Proposition}
\newcommand{\tpsi}{\tilde\psi}
\newcommand{\teps}{\tilde\epsilon}
\newcommand{\Sbar}{\bar S}
\newcommand{\Sigt}{\Sigma_t}
\newcommand{\E}{\mathbb{E}}

\title{Residual Monte Carlo with a continuous linear spatial trial space}
\date{}

\begin{document}
\maketitle

\noindent Slab geometry along $z$ with cells $k=0..K-1$ on nodes $z_0<\dots<z_K$,
widths $h_k$. The transport operator is
$L\psi = \mu\,\partial_z\psi + \Sigt\psi - \mathcal S\psi$, with $\mathcal S$ the
scattering and fission operator (local in $z$), and RMC iterates
\begin{equation}
  r^{(n)} = q - L\tpsi^{(n)},\qquad L\epsilon^{(n)} = r^{(n)},\qquad
  \tpsi^{(n+1)} = \tpsi^{(n)} + \Pi\epsilon^{(n)},
\end{equation}
solving the middle equation by Monte Carlo with a source sampled from $r^{(n)}$ (signed
weights) and $\Pi$ the projection onto the trial space [1--4, 10]. The energy-angle part
of the trial space is any of the discontinuous spaces of
\texttt{linear\_energy\_basis.tex} and \texttt{linear\_discontinuous\_angle.tex}; it
must be the same in every cell (Section~7).

\section{Why: the face-residual pitfall}

With a piecewise-constant $\tpsi$ in $z$, $\mu\,\partial_z\tpsi$ is a sum of delta
functions at the faces, and the residual carries the face sources
\begin{equation}
  r_e = -\mu\,\big(\tpsi_{k}-\tpsi_{k-1}\big)\,\delta(z-z_k)
\end{equation}
(the edge term of \texttt{rmc\_math.tex}). For a flux that varies smoothly in $z$, each
face carries a source $O(h\,\partial_z\psi)$, which does not vanish as the iteration
converges, and the response to it nearly cancels the response to the collision residual
of the adjacent cells: the iteration samples large positive and negative contributions
whose sum is small. Vermaak and Morel identified this as the main pitfall of RMC [1]:
on a 1D slab, a piecewise-constant approximation made RMC $6.5\times$ \emph{less}
efficient than standard Monte Carlo, while a continuous piecewise-linear one made it
$3\times$ more efficient in their Sec.~4.1 benchmark; these factors are not general
performance guarantees. Continuous-energy RMC with piecewise-constant bins shows the same
plateau in 1D [11] (the fuel rod of [D], Fig.~4.9; \texttt{NOTES.md} \S2b). Linear \emph{discontinuous} trial spaces (the classical discontinuous
finite elements of deterministic transport [5, 9]), used by the exponentially
convergent Monte Carlo (ECMC) work [2--4], reduce the jumps to
$O(h^2)$ for smooth fluxes but keep a face source on every face.

In the absence of internal surface sources, the true angular flux is continuous
in $z$ for $\mu\neq0$, also across material
interfaces (it is continuous along characteristics; only its derivative jumps). A
continuous trial space is therefore consistent with the solution, and it removes the
interior face residual altogether.

\section{Trial space and boundary conditions}

With the hat functions $h_i(z)$ ($h_i(z_l)=\delta_{il}$, linear on each cell),
\begin{equation}
  \tpsi(z,E,\mu) = \sum_{i=0}^{K} h_i(z)\,\Phi_i(E,\mu),\qquad
  \tpsi = (1-s)\,\Phi_k + s\,\Phi_{k+1}\ \text{ on cell } k,\quad s=\frac{z-z_k}{h_k},
\end{equation}
where each nodal function $\Phi_i$ lies in the energy-angle trial space (coefficients
$\Phi_{i,gja}$). The space has $K+1$ nodal unknowns per energy-angle coefficient
(linear discontinuous: $2K$; piecewise constant: $K$).

\paragraph{Strong boundary conditions.} The exact solution satisfies the boundary
conditions on incoming directions exactly, so they are imposed on the trial space:
\begin{itemize}
  \item vacuum at $z_0$: $\Phi_{0,gj\cdot}=0$ for bins with $\mu>0$ (and at $z_K$ for
  $\mu<0$); the bins must not straddle $\mu=0$, which the current grids already require;
  \item reflective at $z_0$: $\Phi_0(E,\mu)=\Phi_0(E,-\mu)$, i.e.\ the coefficients of a
  bin and of its mirror bin are one unknown (with the sign of odd angular coefficients
  flipped, see the angle note).
\end{itemize}
With these constraints the trial space contains only functions that satisfy the
boundary conditions, and the boundary face residual of the piecewise-constant scheme
also disappears. This requires an angular space that can impose incoming values
without constraining outgoing ones. Otherwise boundary residuals remain; in
particular this restriction matters for global isotropic or $P_1$ angle in 3D.

\section{Projection and the fixed point}

Let $\Pi$ be the $L^2$-orthogonal projection in $z$ onto the constrained space (tensor
with the energy-angle projection). For each energy-angle coefficient it solves
\begin{equation}
  \sum_l \mathsf M_{il}\,c_l = t_i,\qquad
  \mathsf M_{il} = \int h_ih_l\dd z,\qquad t_i = \int \epsilon\,h_i\dd z,
  \label{eq:mass}
\end{equation}
with $\mathsf M$ tridiagonal (each cell adds $h_k/3$ to both endpoint diagonal
entries and $h_k/6$ to both symmetric off-diagonal entries), and the constrained
unknowns removed or merged. The mass matrix is symmetric positive definite;
conditioning depends on mesh grading and can require scaling. Its direct solve costs
$O(K)$ per coefficient.

\begin{proposition}[Unbiased full-residual update to the constrained $L^2$ projection]
If $L^{-1}$ exists, the full residual (including source and boundary remainders) is
evaluated exactly, and its transport tally is unbiased with finite expectation, then
$\E[\teps\mid\tpsi]=\Pi L^{-1}r=\Pi\psi-\tpsi$ and
$\E[\tpsi^{(n+1)}\mid\tpsi]=\Pi\psi$. This is the best approximation in the
constrained trial space in $L^2$; it is not a proof of convergence of stochastic
iterates or of their variances.
\end{proposition}

\noindent This is the property of \texttt{rmc\_math.tex} (Prop.~1), with $\Pi$ now a
global projection. It matters here because a Galerkin method with continuous trial
\emph{and} test functions applied to the first-order transport operator can exhibit
oscillations and require stabilization in convection-dominated regimes [5, 6].
The full-residual RMC update does not solve those Galerkin
equations: the solve is the exact Monte Carlo transport of the residual, and only the
projection is continuous. The expectation identity needs no Galerkin stabilization,
but sampling and mean-square convergence still need their own analysis. The projection itself can
over- and undershoot near steep boundary layers (as any $L^2$ projection onto
continuous linears does), which is resolved by local refinement, not by changing the
operator.

The collision-only phase uses a modified residual $q_h-L_h\tpsi$, so its fixed
point instead satisfies $\Pi L^{-1}(q_h-L_h\tpsi^*)=0$. Energy and spatial
transport are coupled: exact spatial interpolation alone does not make that fixed
point $\Pi\psi$ for the original physical problem. The full-residual identity
above and the phase-1 contraction conditions in \texttt{rmc\_math.tex} must be
kept separate.

Strong boundary constraints can also remove constants from the spatial trial
space. Its $L^2$ projection then need not preserve the integral of the exact
flux. Projected nodal coefficients and direct physical integral tallies are
different quantities and should be compared with matching reference observables.

\section{Estimating the projection}

The moments $t_i$ are track-length integrals. Inside one cell a track segment is
straight, so $x = 2s-1$ is linear along it and
\begin{equation}
  \int_{\text{segment}} h(z)\dd\ell = \ell\,\frac{h(z_{\rm in})+h(z_{\rm out})}{2}
\end{equation}
exactly for any $h$ linear in $z$. The estimator of $t_i$ therefore only needs, per
cell, the zeroth and first spatial Legendre moments of the track-length estimator,
\begin{equation}
  \hat m_{k,0} = \frac1N\sum w\,\ell,\qquad
  \hat m_{k,1} = \frac1N\sum w\,\ell\,\frac{x_{\rm in}+x_{\rm out}}{2},
\end{equation}
assembled as $\hat t_k \mathrel{+}= (\hat m_{k,0}-\hat m_{k,1})/2$,
$\hat t_{k+1}\mathrel{+}=(\hat m_{k,0}+\hat m_{k,1})/2$ (since $1-s=(1-x)/2$ and
$s=(1+x)/2$). These are unbiased track-length estimators of Legendre moments, as in
functional expansion tallies [7] and linear discontinuous implicit Monte Carlo [8].
Although $\E[x^2]=1/3$ under uniform point sampling, this does not establish a
fixed relative-noise ratio for track-length moments. Weights, track lengths and
within-history correlations determine their covariance; the relative error of
a zero slope is undefined. The required histories must be assessed for the
actual coefficients and observables [7].

\section{Residual}

On cell $k$ (material $m_k$), with $\Delta_k = (\Phi_{k+1}-\Phi_k)/h_k$,
\begin{equation}
  r(z,E,\mu) = \underbrace{Q + \Sbar(z) - \Sigt(E)\,\tpsi(z) - \mu\,\Delta_k}_{r_c}
  \;+\; \underbrace{\mathcal S\tpsi - \Sbar}_{r_s} + (q-Q),
  \qquad
  \Sbar(z) = (1-s)\,\Sbar_{m_k}[\Phi_k] + s\,\Sbar_{m_k}[\Phi_{k+1}],
\end{equation}
\begin{itemize}
  \item \emph{No projection in $z$ is needed for the scattering source}: $\mathcal S$
  acts on energy and angle at each $z$, and $\tpsi$ is linear in $z$ on the cell, so
  $\mathcal S\tpsi$ is exactly the linear interpolation of the nodal sources computed
  with the cell's material moments. The transfer moments are unchanged.
  \item \emph{The streaming term} $-\mu\Delta_k$ replaces the face deltas: it is spread
  over the cell, constant in $z$, linear in $\mu$. As $\tpsi\to\Pi\psi$, $\Delta_k$
  approximates $\partial_z\psi$ and is balanced by the collision terms; what remains is
  the local truncation, without equal-and-opposite pairs.
  \item \emph{No face residual}: interior faces by continuity, boundary faces by the
  strong boundary conditions when those can be imposed in the chosen angular space.
\end{itemize}

\paragraph{Sampling $r_c$.} With piecewise-constant angle, on a bin $(k,g,j)$,
$r_c$ is affine in $(s,\mu)$ for each $E$ (and piecewise quadratic in $E$ for
the linear energy basis). Linear angle adds mixed and quadratic angular terms,
as described in the angle note. With the bin chosen
with probability $A_{kgj}/A$ and $(z,E,\mu)$ uniform in it, the weight
\begin{equation}
  w = \frac{N}{N_c}\,\frac{r_c(z,E,\mu)\,h_k\Delta E_g\Delta\mu_j\,A}{A_{kgj}}
\end{equation}
is an unbiased estimator of $r_c$ for any masses $A_{kgj}>0$ wherever $r_c$ can be
nonzero. The choice $A_{kgj}=\int_{kgj}|r_c|$ is an $L^1$ mass heuristic.
For uniform sampling within bins, the source-weight second moment is minimized
by probabilities proportional to $\sqrt{V_{kgj}\int_{kgj}r_c^2}$, where
$V_{kgj}=h_k\Delta E_g\Delta\mu_j$; transported tallies additionally depend on
their conditional response. For piecewise-constant angle and fixed $E$, the
$L^1$ integral of the affine function over $(s,\mu)$ has a closed form, leaving
quadrature in $E$.

\section{Transfer moments and material interfaces}

The transfer moments $M_{m}$ act on energy and angle only and are those of the
piecewise-constant scheme. Material interfaces sit at nodes: the nodal value
$\Phi_i$ is shared, and each adjacent cell applies its own operator ($\Sigt$,
$M_m$) on its side, so $\tpsi$ is continuous while individual collision and
streaming terms can jump. For the exact solution these terms balance to satisfy
$L\psi=q$ on each side of the interface.

\section{Compatibility with the energy and angle trial spaces}

Continuity is a property of the reconstructed flux, not of the coefficients. It holds
when every cell uses the same energy-angle functions: piecewise-constant or linear
energy bins, shared discontiguous energy elements, linear discontinuous angle. It fails
for a material-dependent energy weight ($1/\Sigma_{t,m}$,
\texttt{self\_shielded\_basis.tex}): across a material interface the weighted shapes
differ, so a face residual reappears there. If that basis is used, the trial space must
generally admit jumps at material interfaces, unless a common reconstructed
trace is represented explicitly by both adjacent trial spaces.

\section{2D and 3D}

The same construction uses continuous $P_1$ elements on triangles and tetrahedra (or
$Q_1$ on quadrilaterals and hexahedra): nodal functions $\Phi_i$, a sparse symmetric
positive-definite mass matrix (a sparse solve, with preconditioning as needed;
the consistent mass, not a lumped one, keeps $\Pi$ a projection), and track-length
tallies of the shape functions, exact as the segment length times the mean of the
endpoint values for $P_1$ (and by Gauss quadrature along the segment for $Q_1$). The
streaming residual $-\boldsymbol\Omega\cdot\nabla\tpsi$ is constant per $P_1$ element
and linear in $\boldsymbol\Omega$; there is no interior face residual.
Incoming conditions can be imposed strongly on boundary nodes only when the
angular space resolves the face's incoming and outgoing directions independently.
For a global isotropic or $P_1$ function $a+\mathbf b\cdot\boldsymbol\Omega$,
vanishing on an incoming hemisphere forces it to vanish on the whole sphere,
incorrectly constraining the outgoing trace. Retain the boundary residual for
these spaces, as in [1], or use an appropriate half-range representation.

\begin{thebibliography}{99}
\bibitem{vermaak} J.~I.~C.~Vermaak and J.~E.~Morel, ``Transport Error Estimation Using
  Residual Monte Carlo,'' \emph{J.\ Comput.\ Phys.} (2022);
  https://www.sciencedirect.com/science/article/pii/S0021999122003680, OSTI 1977284.
\bibitem{peterson} J.~R.~Peterson, J.~E.~Morel and J.~C.~Ragusa,
  ``Exponentially-convergent Monte Carlo for the 1-D transport equation,'' M\&C 2013;
  J.~R.~Peterson, M.S.\ thesis, Texas A\&M University (May 2014),
  \href{https://oaktrust.library.tamu.edu/bitstreams/69ab2dae-3fef-4daf-96d0-2bd1896d81e3/download}
  {university copy}.
  Published formulation: \emph{Residual Monte Carlo for the One-Dimensional
  Particle Transport Equation}, SIAM J.\ Sci.\ Comput.\ \textbf{38}(6),
  B941--B961 (2016), \href{https://doi.org/10.1137/16M1057619}{doi:10.1137/16M1057619}.
\bibitem{franke} B.~C.~Franke, D.~E.~Bruss and J.~E.~Morel, ``Residual Monte Carlo with
  Discrete Scattering Angles in the 1-D Transport Equation,'' ANS Winter Meeting
  \emph{Trans.\ Am.\ Nucl.\ Soc.} \textbf{109}(1), 679--682 (2013),
  \href{https://www.ans.org/pubs/transactions/volume-109/}{ANS publication record}.
\bibitem{bolding} S.~R.~Bolding and J.~E.~Morel, ``Residual Monte Carlo Transport in
  Time with Consistent Low-Order Acceleration for 1D Thermal Radiative Transfer,''
  M\&C 2017,
  \href{https://www.kns.org/files/int_paper/paper/MC2017_2017_2/P102S02-03BoldingS.pdf}
  {conference paper}.
\bibitem{lewis} E.~E.~Lewis and W.~F.~Miller, \emph{Computational Methods of Neutron
  Transport}, Wiley (1984) (finite-element spatial discretizations, linear
  discontinuous methods).
\bibitem{brooks} A.~N.~Brooks and T.~J.~R.~Hughes, ``Streamline upwind/Petrov--Galerkin
  formulations for convection dominated flows,'' \emph{Comput.\ Methods Appl.\ Mech.\
  Eng.} \textbf{32}, 199 (1982).
\bibitem{griesheimer} D.~P.~Griesheimer, W.~R.~Martin and J.~P.~Holloway,
  ``Convergence properties of Monte Carlo functional expansion tallies,''
  \emph{J.\ Comput.\ Phys.} \textbf{211}, 129 (2006).
\bibitem{wollaeger} R.~T.~Wollaeger and A.~B.~Wollaber, ``Implicit Monte Carlo with a
  Linear Discontinuous Finite Element Material Solution and Piecewise Non-Constant
  Opacity,'' LA-UR-15-24303, OSTI 1255243 (2016).
\bibitem{reed} W.~H.~Reed and T.~R.~Hill, ``Triangular mesh methods for the neutron
  transport equation,'' LA-UR-73-479, Los Alamos (1973); P.~Lesaint and P.~A.~Raviart,
  ``On a finite element method for solving the neutron transport equation,'' in
  \emph{Mathematical Aspects of Finite Elements in Partial Differential Equations},
  Academic Press (1974).
\bibitem{halton} J.~H.~Halton, ``Sequential Monte Carlo,'' \emph{Proc.\ Cambridge
  Philos.\ Soc.} \textbf{58}, 57 (1962).
\bibitem{larsen} M.~Larsen et al., ``A Residual Monte Carlo Algorithm for Continuous
  Energy Neutron Transport with Elastic Scattering,'' \emph{Nucl.\ Sci.\ Eng.}
  \textbf{200}(3), 525--538 (2026), doi:10.1080/00295639.2025.2495608;
  ``A Residual Monte Carlo Algorithm for 1D Continuous-Energy
  Neutron Transport,'' \emph{Nucl.\ Sci.\ Eng.} (2026), doi:10.1080/00295639.2026.2659992;
  M.~Larsen, Ph.D.\ dissertation [D], Oregon State University (2025).
\end{thebibliography}

\end{document}
